\begin{abstract}
Large Language Model (LLM)-based Document-level Relation Extraction (DocRE) now matches or exceeds encoder baselines in extraction accuracy, yet outputs bare relation triplets without identifying supporting sentences. Encoder-based evidence methods (SAIS, DREEAM) provide evidence but trail LLM approaches in accuracy, creating a gap between extraction performance and verifiability. We study evidence-grounded DocRE as structured generation and compare Group Relative Policy Optimization (GRPO), Direct Preference Optimization (DPO), and Rejection Sampling Fine-Tuning (RSFT) across three fully evaluated scales (1.7B, 4B, 8B) with 14B validation, two model families, and two datasets. The evidence tax---the extraction cost of generating evidence---is scale-dependent: beneficial at 1.7B ($-$5.63pp Ign-F1) but approximately neutral at $\geq$4B, challenging the assumption that evidence always harms accuracy. Gradient-based optimization amplifies evidence-extraction competition (GRPO catastrophic, DPO gradual), while RSFT avoids decoupling because discrete output selection structurally cannot exploit the extraction-evidence trade-off. We define the Evidence-Decoupling Ratio (EDCR) to quantify this misalignment, validated at per-relation level ($\rho{=}{-}0.79$, $p{<}0.001$). Our 1.7B RSFT model approaches 8B supervised fine-tuning (SFT) performance at 4.7$\times$ fewer parameters, and test-set evaluation preserves model ranking across all scales\footnote{Code and data are available at \url{https://anonymous.4open.science/r/faith-evidence-ie-18BD/}.}.
\end{abstract}
